% \thispagestyle{empty}
% \begin{center}
%     \Large \textbf{Abstract}
% \end{center}


\chapter*{Sommario}
\markboth{Sommario}{Sommario}
\noindent
La terapia fagica si sta affermando come una promettente alternativa agli antibiotici contro i batteri resistenti ai 
farmaci; tuttavia, la sua applicazione clinica dipende dalla corretta identificazione dell'ospite batterico suscettibile 
all'infezione da parte del fago. Gli attuali strumenti computazionali per la predizione dell'ospite forniscono stime 
puntuali o punteggi scalari non calibrati, privi di garanzie a campione finito sulla presenza dell'effettivo ospite tra 
i candidati selezionati.

Questo lavoro propone un framework di predizione conforme (conformal prediction) distribution-free per la classificazione 
fago-ospite atto a colmare tale lacuna. Partendo dall'approccio basato su Conformal Set Aggregation (CSA), ne abbiamo 
sostituito gli inviluppi convessi di calibrazione — suscettibili di un'eccessiva espansione in presenza di cluster di 
punteggi ben separati — introducendo tre geometrie alternative stellate e non convesse (Radiale, a Striscia e 1D 
Collassata), adatte a compiti di classificazione con regioni disgiunte. Per ciascuna configurazione geometrica è 
stata ricavata una regola di calibrazione in forma chiusa che sostituisce la ricerca binaria iterativa della CSA 
standard con un singolo passaggio di ordinamento. Abbiamo inoltre integrato una procedura di calibrazione condizionata 
alla classe, garantendo la validità a campione finito subordinata all'ipotesi di scambiabilità all'interno di ciascuna 
classe di ospiti, anziché una generica garanzia marginale. Infine, abbiamo strutturato una procedura a cascata in due 
stadi (Gram-tipo verso ospite) provvista di un limite dimostrato sulla probabilità di mancata copertura congiunta.

Il framework è stato validato empiricamente su un dataset reale di profili funzionali di proteine fagiche, confermando 
una copertura empirica affidabile al livello di confidenza nominale sia su un task binario di classificazione del tipo 
di Gram, sia su un task multiclasse esteso a 54 ospiti candidati. Il pre-filtraggio basato sulla classificazione di 
Gram ha consentito una riduzione delle dimensioni medie dei set di predizione superiore al $60\%$ rispetto alla baseline 
non filtrata, preservando integralmente il livello di copertura.



\chapter*{Abstract}
\markboth{Abstract}{Abstract}
\noindent
Bacteriophage therapy is emerging as a promising alternative to antibiotics against multi drug resistant bacteria,
but its clinical use depends on correctly matching a phage to a bacterial host it can infect. Existing computational
host prediction tools return point predictions or uncalibrated scalar scores, with no finite sample guarantee that
the true host is among the candidates returned.

This report develops a distribution free conformal prediction framework for phage host classification that closes
this gap. We build on the Conformal Set Aggregation (CSA) approach and replaced its convex calibration envelopes,
which may become uninformatively large when the class conditional score clusters are well separated. We used three star shaped,
non convex alternatives (Radial, Strip, and Collapsed 1D) which are suited to the disjoint geometry of classification
tasks. For each geometry we derived a closed form calibration rule that replaces the iterative binary search used in
standard CSA with a single sorting pass. We also introduce a class conditional calibration, which provides class conditional 
finite sample validity under the exchangeability assumptions within each host class, rather than a single marginal one. 
Finally, we use a two stage Gram type, to host cascade with a proven bound on compounded miscoverage.

We validated the framework on a real dataset of bacteriophage protein functional scores, and achieved reliable
empirical coverage at the target confidence level on both a binary Gram type task and a multiclass host
classification task spanning 54 candidate hosts. Gram type pre filtering also reduced average prediction set sizes
by over $60\%$ compared to an unfiltered baseline, and with no loss of coverage.

\vspace{0.3cm}






